% Transcription of Cayley's Collected Mathematical Papers, Volume IV
% PDF pages 526--545 (printed pages 504--523)
% Papers: 295 "On the Construction of the Ninth Point of Intersection..." (conclusion)
%         296 "On the Conics which pass through the Four Foci of a given Conic"
%         297 "On Some Formulae relating to the Distances of a Point from the Vertices of a Triangle, and to the Problem of Tactions"
%         298 "Report on the Progress of the Solution of certain Special Problems of Dynamics" (opening through Art. 38)

\documentclass[11pt]{article}
\usepackage[a4paper,margin=0.65in]{geometry}
\usepackage[T1]{fontenc}
\usepackage{amsmath,amssymb,amsfonts}
\usepackage{graphicx}
\usepackage{pdfpages}
\usepackage[english,shorthands=off]{babel}
\usepackage{fancyhdr}
\usepackage[bookmarks=false,colorlinks=false]{hyperref}
\usepackage[protrusion=true,expansion=false]{microtype}
\emergencystretch=3em
\tolerance=600
\providecommand{\modd}{\mathrm{mod}}
\providecommand{\Disc}{\mathrm{Disc}}
\providecommand{\canont}{}
\providecommand{\asterism}{*}
\providecommand{\Hessian}{H}
\providecommand{\tome}{}
\providecommand{\page}{p.}
\pagestyle{fancy}\fancyhf{}\fancyhead[L]{Pages 526--545}\fancyhead[R]{Cayley --- Collected Papers, Vol. IV}\renewcommand{\headrulewidth}{0.4pt}
\setlength{\headheight}{15pt}

\begin{document}

% ============================================================
% PDF p.526 = printed p.504 (Paper 295 conclusion)
% ============================================================
\begin{center}
504 \hfill ON THE CONSTRUCTION OF THE NINTH POINT OF INTERSECTION \&c. \hfill [295
\end{center}

Next, if $789=0$, we have identically
\[ 789\,\omega 41 = 789\cdot 741\cdot \omega 81\cdot \omega 94 + 781\cdot 794\cdot \omega 89\cdot \omega 41, \]
which when $789=0$ gives
\[ 789\,\omega 41 = 781\cdot 794\cdot \omega 89\cdot \omega 41, \]
and in like manner,
\[ 789\,\omega 42 = 782\cdot 794\cdot \omega 89\cdot \omega 42, \]
\[ 789\,\omega 43 = 783\cdot 794\cdot \omega 89\cdot \omega 43, \]
in which three equations 4 may be changed into 5 and 6 successively. The equation $\Box = 0$ thus becomes
\[ \left| \begin{array}{ccc} 423\cdot \omega 41, & 431\cdot \omega 42, & 412\cdot \omega 43 \\ 523\cdot \omega 51, & 531\cdot \omega 52, & 512\cdot \omega 53 \\ 623\cdot \omega 61, & 631\cdot \omega 62, & 612\cdot \omega 63 \end{array} \right| = 0, \]
or, what is the same thing,
\[ \left| \begin{array}{ccc} 423\cdot 41\omega, & 421\cdot 4\omega 3, & 42\omega\cdot 431 \\ 523\cdot 51\omega, & 521\cdot 5\omega 3, & 52\omega\cdot 531 \\ 623\cdot 61\omega, & 621\cdot 6\omega 3, & 62\omega\cdot 631 \end{array} \right| = 0, \]
and since the sum of the terms in each line of the determinant is $=0$, the determinant is as it should be $=0$.

The foregoing equation
\[ \left| \begin{array}{ccc} 412\cdot 789\,\omega 43, & 423\cdot 789\,\omega 41, & 431\cdot 789\,\omega 42 \\ 512\cdot 789\,\omega 53, & 523\cdot 789\,\omega 51, & 531\cdot 789\,\omega 52 \\ 612\cdot 789\,\omega 63, & 623\cdot 789\,\omega 61, & 631\cdot 789\,\omega 62 \end{array} \right| \]
\[ = e \cdot 123\cdot 789\cdot 78\omega\cdot 79\omega\cdot 89\omega\cdot 123456789\omega, \]
(since $412$, $789\,\omega 43$, \&c. are interpretable functions of the coordinates) affords a geometrical interpretation of the equation
\[ 123456789\omega = 0 \]
between the coordinates of the ten points of the cubic; but it would be more satisfactory if a similar identical equation could be found, having on the right-hand side the function $123456789\omega$ without the irrelevant factor
\[ 123\cdot 789\cdot 78\omega\cdot 79\omega\cdot 89\omega. \]

\medskip
\noindent\emph{2, Stone Buildings, W.C., 6th March, 1862.}

\vfill
\begin{flushright}\small 296]\quad 505\end{flushright}

% ============================================================
% PDF p.527 = printed p.505 (Paper 296 opening)
% ============================================================
\newpage
\begin{center}
\textbf{\large 296.}\\[1ex]
\textbf{ON THE CONICS WHICH PASS THROUGH THE FOUR FOCI\\ OF A GIVEN CONIC.}
\end{center}

\begin{center}\small
[From the \emph{Quarterly Journal of Pure and Applied Mathematics}, vol.\ v.\ (1862), pp.\ 275--280.]
\end{center}

\medskip

The foci of a conic are the points of intersection of the tangents through the circular points at infinity; the pair of tangents through each of the circular points at infinity is a conic through the four foci; and we have thus two conics $P=0$, $Q=0$ passing through the four foci; the equation of any other conic through the four foci is of course $P+\lambda Q=0$; and in particular if $\lambda$ be suitably determined this equation gives the axes of the conic.

I was led to develope the solution, in seeking to obtain the elegant formulae given in Mr P.\ J.\ Hensley's paper ``Determination of the foci of the conic section expressed by trilinear coordinates,'' \emph{Journal}, t.\ v., pp.\ 177--183, (March, 1862).

I take the coordinates to be proportional to the perpendicular distances of the point from the sides of the fundamental triangle, each distance divided by the perpendicular distance of the side from the opposite angle. This being so, the equation of the line infinity is
\[ \frac{\alpha}{x}+\frac{\beta}{y}+\frac{\gamma}{z}=0, \]
and, $\alpha,\beta,\gamma$ denoting the sides of the fundamental triangle, the equation of the circle circumscribed about the triangle is
\[ \frac{\alpha^{2}}{x}+\frac{\beta^{2}}{y}+\frac{\gamma^{2}}{z}=0. \]
The foregoing two equations determine the circular points at infinity; and if $(x_{1},y_{1},z_{1})$ are the coordinates, there is no difficulty in obtaining the system of values
\[
\begin{aligned}
x_{1} : y_{1} : z_{1} &= -\alpha \,:\, \beta(\cos C + i\sin C) \,:\, \gamma(\cos B + i\sin B)\\
&= \alpha(\cos C - i\sin C) \,:\, -\beta \,:\, \gamma(\cos A + i\sin A)\\
&= \alpha(\cos B + i\sin B) \,:\, \beta(\cos A - i\sin A) \,:\, -\gamma,
\end{aligned}
\]

\vfill
\begin{flushright}\small C.\ IV.\quad 64\end{flushright}

% ============================================================
% PDF p.528 = printed p.506
% ============================================================
\newpage
\begin{center}
506 \hfill ON THE CONICS WHICH PASS THROUGH \hfill [296
\end{center}

where as usual $i=\sqrt{(-1)}$, and where $A$, $B$, $C$ denote the angles of the triangle, so that the cosines and sines of these angles denote given functions of $\alpha,\beta,\gamma$. The coordinates $(x_{2},y_{2},z_{2})$ of the other circular point at infinity are of course obtained by merely writing $-i$ for $i$. We find also
\[ x_{1}x_{2} : y_{1}y_{2} : z_{1}z_{2} : y_{1}z_{2}+y_{2}z_{1} : z_{1}x_{2}+z_{2}x_{1} : x_{1}y_{2}+x_{2}y_{1} \]
\[ = -\alpha^{2} : -\beta^{2} : -\gamma^{2} : \beta^{2}+\gamma^{2}-\alpha^{2} : \gamma^{2}+\alpha^{2}-\beta^{2} : \alpha^{2}+\beta^{2}-\gamma^{2}, \]
which are the formulae chiefly made use of in the sequel.

Suppose now that the equation of the conic is
\[ U=(a,b,c,f,g,h \mid x,y,z)^{2}=0; \]
then putting for a moment
\[ U_{1}=(a,b,c,f,g,h \mid x_{1},y_{1},z_{1})^{2}, \]
\[ V_{1}=(a,b,c,f,g,h \mid x,y,z\mid x_{1},y_{1},z_{1}), \]
and the like as regards $U_{2}$ and $W_{2}$; the tangents from the points $(x_{1},y_{1},z_{1})$ and $(x_{2},y_{2},z_{2})$ respectively are
\[ U\,U_{1}-V_{1}^{2}=0, \]
\[ U\,U_{2}-W_{2}^{2}=0, \]
which are respectively pairs of lines intersecting in the four foci. And it is moreover clear that the equation of the axes is
\[ U_{1}W_{2}^{2}-U_{2}V_{1}^{2}=0. \]
The foregoing equations may be written
\[ (\mathfrak{A},\mathfrak{B},\mathfrak{C},\mathfrak{F},\mathfrak{G},\mathfrak{H} \mid yz_{1}-y_{1}z,\; zx_{1}-z_{1}x,\; xy_{1}-x_{1}y)^{2}=0, \]
\[ (\mathfrak{A},\mathfrak{B},\mathfrak{C},\mathfrak{F},\mathfrak{G},\mathfrak{H} \mid yz_{2}-y_{2}z,\; zx_{2}-z_{2}x,\; xy_{2}-x_{2}y)^{2}=0, \]
where $(\mathfrak{A},\mathfrak{B},\mathfrak{C},\mathfrak{F},\mathfrak{G},\mathfrak{H})$ are the inverse system of coefficients.

These may be written
\[ (\mathfrak{a},\mathfrak{b},\mathfrak{c},\mathfrak{f},\mathfrak{g},\mathfrak{h} \mid x_{1},y_{1},z_{1})^{2}=0, \]
\[ (\mathfrak{a},\mathfrak{b},\mathfrak{c},\mathfrak{f},\mathfrak{g},\mathfrak{h} \mid x_{2},y_{2},z_{2})^{2}=0, \]
where $\mathfrak{a},\mathfrak{b},\mathfrak{c},\mathfrak{f},\mathfrak{g},\mathfrak{h}$ are quadric functions of $x,y,z$, viz.
\[ \begin{aligned}
\mathfrak{a} &= \mathfrak{B}z^{2}+\mathfrak{C}y^{2}-2\mathfrak{F}yz,\\
\mathfrak{b} &= \mathfrak{C}x^{2}+\mathfrak{A}z^{2}-2\mathfrak{G}zx,\\
\mathfrak{c} &= \mathfrak{A}y^{2}+\mathfrak{B}x^{2}-2\mathfrak{H}xy,\\
\mathfrak{f} &= -\mathfrak{A}yz-\mathfrak{F}x^{2}+\mathfrak{G}xy+\mathfrak{H}xz,\\
\mathfrak{g} &= -\mathfrak{B}zx-\mathfrak{G}y^{2}+\mathfrak{H}yz+\mathfrak{F}yx,\\
\mathfrak{h} &= -\mathfrak{C}xy-\mathfrak{H}z^{2}+\mathfrak{F}zx+\mathfrak{G}zy,
\end{aligned} \]

% ============================================================
% PDF p.529 = printed p.507
% ============================================================
\newpage
\begin{center}
296] \hfill THE FOUR FOCI OF A GIVEN CONIC. \hfill 507
\end{center}

and this being so, I combine the two equations as follows:
\[ \begin{aligned}
x_{2}^{2}(\mathfrak{a},\ldots\mid x_{1},y_{1},z_{1})^{2}+x_{1}^{2}(\mathfrak{a},\ldots\mid x_{2},y_{2},z_{2})^{2} &=0,\\
y_{2}^{2}\ \text{''}\ \phantom{(\mathfrak{a},\ldots\mid x_{1},y_{1},z_{1})^{2}}+y_{1}^{2}\ \text{''} &=0,\\
z_{2}^{2}\ \text{''}\ \phantom{(\mathfrak{a},\ldots\mid x_{1},y_{1},z_{1})^{2}}+z_{1}^{2}\ \text{''} &=0,\\
y_{2}z_{2}(\mathfrak{a},\ldots\mid x_{1},y_{1},z_{1})^{2}+y_{1}z_{1}(\mathfrak{a},\ldots\mid x_{2},y_{2},z_{2})^{2} &=0,\\
z_{2}x_{2}\ \text{''}\ \phantom{(\mathfrak{a},\ldots\mid x_{1},y_{1},z_{1})^{2}}+z_{1}x_{1}\ \text{''} &=0,\\
x_{2}y_{2}\ \text{''}\ \phantom{(\mathfrak{a},\ldots\mid x_{1},y_{1},z_{1})^{2}}+x_{1}y_{1}\ \text{''} &=0,
\end{aligned} \]
any one of which is the equation of a conic passing through the four foci; the current coordinates being always $(x,y,z)$.

The first of these equations is
\[ \mathfrak{a}(-2x_{1}^{2}x_{2}^{2})+\mathfrak{b}(x_{1}^{2}y_{2}^{2}+x_{2}^{2}y_{1}^{2})+\mathfrak{c}(x_{1}^{2}z_{2}^{2}+x_{2}^{2}z_{1}^{2}) \]
\[ +2\mathfrak{f}(x_{1}^{2}y_{2}z_{2}+x_{2}^{2}y_{1}z_{1})+2\mathfrak{g}(x_{1}^{2}z_{2}x_{1}+x_{1}^{2}z_{1}x_{2})+2\mathfrak{h}(x_{2}^{2}x_{1}y_{1}+x_{1}^{2}x_{2}y_{2})=0, \]
where the quantities multiplied by $\mathfrak{a}$, $\mathfrak{b}$, \&c.\ are all of them easily expressible in terms of $x_{1}x_{2}$, $y_{1}y_{2}$, $z_{1}z_{2}$, $y_{1}z_{2}+y_{2}z_{1}$, $z_{1}x_{2}+z_{2}x_{1}$, $x_{1}y_{2}+x_{2}y_{1}$, which are respectively proportional to given functions of $(\alpha,\beta,\gamma)$; and replacing for $\mathfrak{a}$, $\mathfrak{b}$, \&c.\ their values, the equation is
{\small\begin{align*}
&(\mathfrak{B}z^{2}+\mathfrak{C}y^{2}-2\mathfrak{F}yz)\cdot 2\alpha^{4}\\
&\quad+(\mathfrak{C}x^{2}+\mathfrak{A}z^{2}-2\mathfrak{G}zx)\cdot \{(\alpha^{2}+\beta^{2}-\gamma^{2})^{2}-2\alpha^{2}\beta^{2}\}\\
&\quad+(\mathfrak{A}y^{2}+\mathfrak{B}x^{2}-2\mathfrak{H}xy)\cdot \{(\gamma^{2}+\alpha^{2}-\beta^{2})^{2}-2\gamma^{2}\alpha^{2}\}\\
&\quad+2(-\mathfrak{A}yz-\mathfrak{F}x^{2}+\mathfrak{G}xy+\mathfrak{H}xz)\cdot \alpha^{2}(\beta^{2}+\gamma^{2}-\alpha^{2})\\
&\quad+2(-\mathfrak{B}zx-\mathfrak{G}y^{2}+\mathfrak{H}yz+\mathfrak{F}yz)\cdot -\alpha^{2}(\gamma^{2}+\alpha^{2}-\beta^{2})\\
&\quad+2(-\mathfrak{C}xy-\mathfrak{H}z^{2}+\mathfrak{F}zx+\mathfrak{G}zy)\cdot -\alpha^{2}(\alpha^{2}+\beta^{2}-\gamma^{2})=0.
\end{align*}}
Now putting for shortness
\[ \Box = \alpha^{4}+\beta^{4}+\gamma^{4}-2\beta^{2}\gamma^{2}-2\gamma^{2}\alpha^{2}-2\alpha^{2}\beta^{2}, \]
so that $-\Box$ is equal to sixteen times the square of the area of the fundamental triangle, the coefficient of $x^{2}$ is
\[ = \mathfrak{B}(\Box+2\alpha^{2}\gamma^{2})+\mathfrak{C}(\Box+2\alpha^{2}\beta^{2})-2\mathfrak{H}\{-\Box-\alpha^{2}(\beta^{2}+\gamma^{2}-\alpha^{2})\}, \]
which is $= (\mathfrak{B}+\mathfrak{C}+2\mathfrak{H})\,\Box + 2\alpha^{2}\{\mathfrak{B}\gamma^{2}+\mathfrak{C}\beta^{2}+\mathfrak{H}(\beta^{2}+\gamma^{2}-\alpha^{2})\}.$

And reducing in a similar manner the other coefficients, the equation is
\[ \Box\{(\mathfrak{B}+\mathfrak{C}+2\mathfrak{H})x^{2}+\mathfrak{A}(y+z)^{2}-2(\mathfrak{F}+\mathfrak{G})x(y+z)\}+2\alpha^{2}\Theta=0, \]
where for shortness
{\small\begin{align*}
\Theta = {}& x^{2}\cdot \{\mathfrak{B}\gamma^{2}+\mathfrak{C}\beta^{2}+\mathfrak{H}(\beta^{2}+\gamma^{2}-\alpha^{2})\}\\
&+y^{2}\cdot \{\mathfrak{C}\alpha^{2}+\mathfrak{A}\gamma^{2}+\mathfrak{G}(\gamma^{2}+\alpha^{2}-\beta^{2})\}\\
&+z^{2}\cdot \{\mathfrak{A}\beta^{2}+\mathfrak{B}\alpha^{2}+\mathfrak{F}(\alpha^{2}+\beta^{2}-\gamma^{2})\}\\
&+yz\cdot \{-2\mathfrak{F}\alpha^{2}+\mathfrak{A}(\beta^{2}+\gamma^{2}-\alpha^{2})-\mathfrak{G}(\gamma^{2}+\alpha^{2}-\beta^{2})-\mathfrak{H}(\alpha^{2}+\beta^{2}-\gamma^{2})\}\\
&+zx\cdot \{-2\mathfrak{G}\beta^{2}-\mathfrak{F}(\beta^{2}+\gamma^{2}-\alpha^{2})+\mathfrak{B}(\gamma^{2}+\alpha^{2}-\beta^{2})-\mathfrak{H}(\alpha^{2}+\beta^{2}-\gamma^{2})\}\\
&+xy\cdot \{-2\mathfrak{H}\gamma^{2}-\mathfrak{F}(\beta^{2}+\gamma^{2}-\alpha^{2})-\mathfrak{G}(\gamma^{2}+\alpha^{2}-\beta^{2})+\mathfrak{C}(\alpha^{2}+\beta^{2}-\gamma^{2})\},
\end{align*}}

\begin{flushright}\small 64---2\end{flushright}

% ============================================================
% PDF p.530 = printed p.508
% ============================================================
\newpage
\begin{center}
508 \hfill ON THE CONICS WHICH PASS THROUGH \hfill [296
\end{center}

or, what is the same thing, the equation is
\[ \Box\{(\mathfrak{A}+\mathfrak{B}+\mathfrak{C}+2\mathfrak{F}+2\mathfrak{G}+2\mathfrak{H})x^{2}-2(\mathfrak{A}+\mathfrak{F}+\mathfrak{G})x(x+y+z)+\mathfrak{A}(x+y+z)^{2}\}+2\alpha^{2}\Theta=0, \]
or putting for shortness
\[ \begin{aligned}
\mathfrak{A}+\mathfrak{B}+\mathfrak{C}+2\mathfrak{F}+2\mathfrak{G}+2\mathfrak{H} &= \mathfrak{K},\\
\mathfrak{A}+\mathfrak{F}+\mathfrak{G} &= \mathfrak{L},\\
\mathfrak{F}+\mathfrak{B}+\mathfrak{H} &= \mathfrak{M},\\
\mathfrak{G}+\mathfrak{H}+\mathfrak{C} &= \mathfrak{N},
\end{aligned} \]
so that in fact
\[ \mathfrak{L}+\mathfrak{M}+\mathfrak{N} = \mathfrak{K}, \]
the equation is
\[ \Box\{\mathfrak{K}x^{2}-2\mathfrak{L}x(x+y+z)+\mathfrak{A}(x+y+z)^{2}\}+2\alpha^{2}\Theta=0. \]
The equation with $y_{2}z_{2}$, $y_{1}z_{1}$ is in a similar manner found to be
\[ \resizebox{\textwidth}{!}{$\displaystyle
\mathfrak{a}\{\mathfrak{K}x^{2}-[(\mathfrak{G}+\mathfrak{C})y^{2}-(\mathfrak{H}+\mathfrak{B})z^{2}+(\mathfrak{A}+2\mathfrak{F}+\mathfrak{G}+\mathfrak{H})yz+(-\mathfrak{H}-\mathfrak{F}+\mathfrak{G})zx+(-\mathfrak{G}-\mathfrak{F}+\mathfrak{H})xy]\}
$} \]
or, what is the same thing,
\[ \resizebox{\textwidth}{!}{$\displaystyle
\Box\,[(\mathfrak{A}+\mathfrak{B}+\mathfrak{C}+2\mathfrak{F}+2\mathfrak{G}+2\mathfrak{H})yz-\{(\mathfrak{C}+\mathfrak{G}+\mathfrak{H})y+(\mathfrak{F}+\mathfrak{B}+\mathfrak{H})z\}(x+y+z)+\mathfrak{F}(x+y+z)^{2}]
-(\beta^{2}+\gamma^{2}-\alpha^{2})\Theta=0,
$} \]
or, finally,
\[ \Box\{\mathfrak{K}yz-(\mathfrak{N}y+\mathfrak{M}z)(x+y+z)+\mathfrak{F}(x+y+z)^{2}\}-(\beta^{2}+\gamma^{2}-\alpha^{2})\Theta=0. \]
Hence the entire system of equations is
\[ \begin{aligned}
\Box\{\mathfrak{K}x^{2}-2\mathfrak{L}x(x+y+z)+\mathfrak{A}(x+y+z)^{2}\}+2\alpha^{2}\Theta &= 0,\\
\Box\{\mathfrak{K}y^{2}-2\mathfrak{M}y(x+y+z)+\mathfrak{B}(x+y+z)^{2}\}+2\beta^{2}\Theta &= 0,\\
\Box\{\mathfrak{K}z^{2}-2\mathfrak{N}z(x+y+z)+\mathfrak{C}(x+y+z)^{2}\}+2\gamma^{2}\Theta &= 0,\\
\Box\{\mathfrak{K}yz-(\mathfrak{N}y+\mathfrak{M}z)(x+y+z)+\mathfrak{F}(x+y+z)^{2}\}-(\beta^{2}+\gamma^{2}-\alpha^{2})\Theta &= 0,\\
\Box\{\mathfrak{K}zx-(\mathfrak{L}z+\mathfrak{N}x)(x+y+z)+\mathfrak{G}(x+y+z)^{2}\}-(\gamma^{2}+\alpha^{2}-\beta^{2})\Theta &= 0,\\
\Box\{\mathfrak{K}xy-(\mathfrak{M}x+\mathfrak{L}y)(x+y+z)+\mathfrak{H}(x+y+z)^{2}\}-(\alpha^{2}+\beta^{2}-\gamma^{2})\Theta &= 0,
\end{aligned} \]
which are the equations of six conics, each of them passing through the four foci.

From the first three of these, we have
\[ \frac{1}{\alpha^{2}}\{\mathfrak{K}x^{2}-2\mathfrak{L}x(x+y+z)+\mathfrak{A}(x+y+z)^{2}\} \]
\[ = \frac{1}{\beta^{2}}\{\mathfrak{K}y^{2}-2\mathfrak{M}y(x+y+z)+\mathfrak{B}(x+y+z)^{2}\} \]
\[ = \frac{1}{\gamma^{2}}\{\mathfrak{K}z^{2}-2\mathfrak{N}z(x+y+z)+\mathfrak{C}(x+y+z)^{2}\}, \]

% ============================================================
% PDF p.531 = printed p.509
% ============================================================
\newpage
\begin{center}
296] \hfill THE FOUR FOCI OF A GIVEN CONIC. \hfill 509
\end{center}

which, allowing for the difference of notation, are Mr Hensley's equations: it appears by his investigation that their geometrical signification is as follows; viz.\ if for shortness we denote the equations by
\[ A=B=C, \]
then if we consider the tangents parallel to the $x$-side of the fundamental triangle, and the tangents parallel to the $y$-side of the fundamental triangle, the equation $A=B$ is the locus of a point such that the feet of the perpendiculars let fall from it on the four tangents lie in a circle. And similarly for the equations $A=C$, $B=C$.

If we multiply the six equations by $1,1,1,2,2,2$ and add, we obtain the identical equation $0=0$; if we multiply them by $a,b,c,2f,2g,2h$ and add, then after some easy reductions, we obtain for the equation of a new conic passing through the four foci
\[ \Box\{\mathfrak{K}U+K(x+y+z)^{2}\}+2S\Theta=0, \]
where
\[ U=(a,b,c,f,g,h \mid x,y,z)^{2}, \]
$K$ is the discriminant $abc-af^{2}-bg^{2}-ch^{2}+2fgh$, and
\[ S=(a\alpha^{2}+b\beta^{2}+c\gamma^{2}-f(\beta^{2}+\gamma^{2}-\alpha^{2})-g(\gamma^{2}+\alpha^{2}-\beta^{2})-h(\alpha^{2}+\beta^{2}-\gamma^{2}), \]
or, what is the same thing,
\[ S=(f+a-h-g)\alpha^{2}+(g-h+b-f)\beta^{2}+(h-g-f+c)\gamma^{2}. \]
It would be interesting to ascertain the geometrical signification of the six conics and of the last-mentioned new conic.

\medskip
\noindent\emph{2, Stone Buildings, W.C., March 13th, 1862.}

\vfill
\begin{flushright}\small 510\end{flushright}

% ============================================================
% PDF p.532 = printed p.510 (Paper 297 opening)
% ============================================================
\newpage
\begin{center}
510 \hfill\hfill [297
\end{center}

\bigskip
\begin{center}
\textbf{\large 297.}\\[1ex]
\textbf{ON SOME FORMULAE RELATING TO THE DISTANCES OF A\\ POINT FROM THE VERTICES OF A TRIANGLE, AND TO THE\\ PROBLEM OF TACTIONS.}
\end{center}

\begin{center}\small
[From the \emph{Quarterly Journal of Pure and Applied Mathematics}, vol.\ v.\ (1862), pp.\ 381--384.]
\end{center}

\medskip
The relation between the distances of four points $1,2,3,4$ in a plane is
\[ \left| \begin{array}{ccccc} 0 & 12^{2} & 13^{2} & 14^{2} & 1\\ 21^{2} & 0 & 23^{2} & 24^{2} & 1\\ 31^{2} & 32^{2} & 0 & 34^{2} & 1\\ 41^{2} & 42^{2} & 43^{2} & 0 & 1\\ 1 & 1 & 1 & 1 & 0 \end{array} \right| = 0, \]
where, see my paper ``Note on the value of certain Determinants the terms of which are the squared distances of Points in a plane or in space,'' \emph{Quarterly Journal of Mathematics}, t.\ iii., p.\ 275 (1859), [286], the determinant is
\[ = \Sigma\, 12^{2}\cdot 23^{2}\cdot 34^{2}-\Sigma\, 12^{2}\cdot 34^{2}\cdot 43^{2}-\Sigma\, 12^{2}\cdot 23^{2}\cdot 31^{2}, \]
an identity which subsists without the aid of the relations $12=21$, \&c., and in which the $\Sigma$, $\Sigma$, $\Sigma$ contain 24, 12, and 8 terms respectively.

Writing $23=f$, $31=g$, $12=h$, $14=a$, $24=b$, $34=c$, the determinant is
\[ \begin{aligned}
={}& 2\,\{\,g^{2}h^{2}(b^{2}+c^{2})+h^{2}f^{2}(c^{2}+a^{2})+f^{2}g^{2}(a^{2}+b^{2})\\
&+b^{2}c^{2}(g^{2}+h^{2})+c^{2}a^{2}(h^{2}+f^{2})+a^{2}b^{2}(f^{2}+g^{2})\\
&-a^{2}f^{2}(a^{2}+f^{2})-b^{2}g^{2}(b^{2}+g^{2})-c^{2}h^{2}(c^{2}+h^{2})\\
&-b^{2}c^{2}f^{2}-c^{2}a^{2}g^{2}-a^{2}b^{2}h^{2}-f^{2}g^{2}h^{2}\,\}\\
={}& -2D,
\end{aligned} \]

% ============================================================
% PDF p.533 = printed p.511
% ============================================================
\newpage
\begin{center}
297] \hfill ON SOME FORMULAE RELATING TO THE DISTANCES OF A POINT \&c. \hfill 511
\end{center}

if $D$ denote the function in $\{\ \}$ with the signs reversed. The function $D$ may be expressed in the form
\[ D = \begin{aligned}[t]
&+(a^{2}f^{2}+b^{2}c^{2})(f^{2}-g^{2}-h^{2})\\
&+(b^{2}g^{2}+c^{2}a^{2})(g^{2}-h^{2}-f^{2})\\
&+(c^{2}h^{2}+a^{2}b^{2})(h^{2}-f^{2}-g^{2}),
\end{aligned} \]
and also in the form
\[ \Box = U^{2}+(f+g+h)V, \]
if for shortness
\[ U = a^{2}f+b^{2}g+c^{2}h+fgh, \]
\[ \begin{aligned}
V = {}& (a^{2}f^{2}+b^{2}c^{2})(f-g-h)\\
&+(b^{2}g^{2}+c^{2}a^{2})(g-h-f)\\
&+(c^{2}h^{2}+a^{2}b^{2})(h-f-g);
\end{aligned} \]
and it may be remarked that since $D$ is an even function of $f$, $g$, $h$, we may in this last formula change at pleasure the signs of these quantities; we thus obtain in all four similar forms of the function $D$.

It is clear that considering a triangle, and any point in the plane of the triangle, $f$, $g$, $h$ may be taken to denote the sides of the triangle, and $a$, $b$, $c$ the distances of the point from the vertices: and the equation $D=0$ is the relation connecting the sides and distances.

The equation $f+g+h=0$ denotes that the vertices are \emph{in lined}, and when this equation is satisfied we have
\[ U = a^{2}f+b^{2}g+c^{2}h+fgh = 0, \]
which is in fact, as it is easy to see, the relation connecting the distances of a point from any three points \emph{in lined}.

For $a$, $b$, $c$ write $a+x$, $b+x$, $c+x$; $x$ will be the radius of a circle touching the circles, radii $a$, $b$, $c$, described about the vertices as centres. The equation $D=0$ becomes after all reductions
{\small\begin{align*}
&U^{2}-(f+g+h)V\\
&\quad+x\,[\,4U(af+bg+ch)\\
&\qquad-2(f+g+h)\{(a^{2}f^{2}+bc(b+c))(f-g-h)\\
&\qquad\qquad+(b^{2}g^{2}+ca(c+a))(g-h-f)\\
&\qquad\qquad+(c^{2}h^{2}+ab(a+b))(h-f-g)\}\,]\\
&\quad+x^{2}[\,f^{2}\{-4a^{2}+6a(b+c)-6bc\}\\
&\qquad+g^{2}\{-4b^{2}+6b(c+a)-6ca\}\\
&\qquad+h^{2}\{-4c^{2}+6c(a+b)-6ab\}\,] = 0,
\end{align*}}
which is a quadratic equation only: the two circles thus obtained are those which touch the given circles all three externally or all three internally. But by changing in every possible manner the signs of $a$, $b$, $c$ we obtain in all four equations giving the eight tangent circles. It may be noticed that if as before $f+g+h=0$, $U=0$,

% ============================================================
% PDF p.534 = printed p.512
% ============================================================
\newpage
\begin{center}
512 \hfill ON SOME FORMULAE RELATING TO THE DISTANCES OF A POINT \&c. \hfill [297
\end{center}

then not only the constant term vanishes, but the coefficient of $x$ also vanishes or the equation becomes simply $x^{2}=0$.

In particular, suppose $f=b+c$, $g=c+a$, $h=a+b$; developing this \emph{de novo}, and putting for shortness
\[ \begin{aligned}
a+b+c &= p,\\
bc+ca+ab &= q,\\
abc &= r,
\end{aligned} \]
we find
\[ U = 2\{px^{2}+2qx+pq-2r\}, \]
\[ V = 2\{px^{4}+4qx^{3}+(2pq+12r)x^{2}+4q^{2}x+pq^{2}-4qr\}, \]
and then the equation $D = U^{2}-2pV = 0$ gives
\[ \tfrac{1}{4}D = (px^{2}+2qx+pq-2r)^{2}-p\{px^{4}+4qx^{3}+(2pq+12r)x^{2}+4q^{2}x+pq^{2}-4qr\} \]
\[ = 4\{(q^{2}-4pr)x^{2}-2qrx+r^{2}\}, \]
so that we have
\[ \tfrac{1}{16}D = (q^{2}-4pr)x^{2}-2qrx+r^{2} = (qx-r)^{2}-4prx^{2} = 0, \]
and thence
\[ qx-r = \pm x\sqrt{pr},\quad\text{or}\quad x = \frac{r}{q\mp\sqrt{pr}}, \]
which gives the radii of the circles inscribed in and circumscribed about the three circles radii $a$, $b$, $c$, whereof each touches the two others: a formula given by Descartes, \emph{Epistolae} (Ed.\ 2, Franc.\ 1792), Pars iii., p.\ 261, in a letter to the Princess Elizabeth, viz.\ Descartes has
\[ \resizebox{\textwidth}{!}{$\displaystyle
(d^{2}e^{2}+d^{2}f^{2}+e^{2}f^{2}-2def^{2}-2d^{2}ef-2de^{2}f)x^{2}-2(de^{2}f^{2}+d^{2}ef^{2}+d^{2}e^{2}f)x+d^{2}e^{2}f^{2} = 0,
$} \]
which putting $a$, $b$, $c$ for his $d$, $e$, $f$, becomes \emph{ut supra}
\[ x^{2}(q^{2}-4pr)-2qrx+r^{2} = 0. \]

In conclusion I notice the following formula which is obtained without difficulty, viz.\ if as before we have a triangle the sides whereof are $f$, $g$, $h$, and if $a$, $b$, $c$ are the distances of a point from the vertices (so that as before $D=0$) then the perpendicular distances of the point from the sides, each perpendicular distance divided by the perpendicular distance of the opposite vertex from the same side, are as follows: viz.\ the quotient for the side $f$ is
\[ = \frac{1}{16\Delta^{2}}\,[(b^{2}-c^{2})(g^{2}-h^{2})+f^{2}(b^{2}+c^{2}+g^{2}+h^{2}-2a^{2})-f^{4}], \]
where $\Delta$ is the area of the triangle. It is clear that we ought to have
\[ \Sigma\,[(b^{2}-c^{2})(g^{2}-h^{2})+f^{2}(b^{2}+c^{2}+g^{2}+h^{2}-2a^{2})-f^{4}] = 16\Delta^{2}, \]
and this equation in fact reduces itself to
\[ 2g^{2}h^{2}+2h^{2}f^{2}+2f^{2}g^{2}-f^{4}-g^{4}-h^{4} = 16\Delta^{2}, \]
which is right.

\medskip
\noindent\emph{2, Stone Buildings, W.C., 11th September, 1862.}

% ============================================================
% PDF p.535 = printed p.513 (Paper 298 opening)
% ============================================================
\newpage
\begin{center}
298] \hfill\hfill 513
\end{center}

\bigskip
\begin{center}
\textbf{\large 298.}\\[1ex]
\textbf{REPORT ON THE PROGRESS OF THE SOLUTION OF CERTAIN\\ SPECIAL PROBLEMS OF DYNAMICS.}
\end{center}

\begin{center}\small
[From the \emph{Report of the British Association for the Advancement of Science}, 1862, pp.\ 184--252.]
\end{center}

\medskip

My ``Report on the Recent Progress of Theoretical Dynamics'' was published in the Report of the British Association for the year 1857, [195]. The present Report (which is in some measure supplemental thereto) relates to the Special Problems of Dynamics: to give a general idea of the contents, I will at once mention the heads under which these problems are considered; viz., relating to the motion of a particle or system of particles, we have
\begin{quote}
Rectilinear Motion;\\
Central Forces, and in particular\\
Elliptic Motion;\\
The Problem of two Centres;\\
The Spherical Pendulum;\\
Motion as affected by the Rotation of the Earth, and Relative Motion in general;\\
Miscellaneous Problems;\\
The Problem of three bodies.
\end{quote}

And relating to the motion of a solid body, we have
\begin{quote}
The Transformation of Coordinates;\\
Principal Axes, and Moments of Inertia;\\
Rotation of a Solid Body;\\
Kinematics of a Solid Body;\\
Miscellaneous Problems.
\end{quote}

\vfill
\begin{flushright}\small C.\ IV.\quad 65\end{flushright}

% ============================================================
% PDF p.536 = printed p.514
% ============================================================
\newpage
\begin{center}
514 \hfill REPORT ON THE PROGRESS OF THE SOLUTION OF \hfill [298
\end{center}

As regards the first division of the subject, I remark that the lunar and planetary theories, as usually treated, do not (properly speaking) relate to the problem of three bodies, but to that of disturbed elliptic motion---a problem which is not considered in the present Report. The problem of the spherical pendulum is that of a particle moving on a spherical surface; but, with this exception, I do not much consider the motion of a particle on a given curve or surface, nor the motion in a resisting medium; what is said on these subjects is included under the head Miscellaneous Problems. The first six heads relate exclusively, and the head Miscellaneous Problems relates principally, to the motion of a single particle. As regards the second division of the subject, I will only remark that, from its intimate connexion with the theory of the motion of a solid body, I have been induced to make a separate head of the geometrical subject, ``Transformation of Coordinates,'' and to treat of it in considerable detail.

I have inserted at the end of the present Report a list of the memoirs and works referred to, arranged (not, as in the former Report, in chronological order, but) alphabetically according to the authors' names: those referred to in the former Report formed for the purpose thereof a single series, which is not here the case. The dates specified are for the most part those on the title-page of the volume, being intended to show approximately the date of the researches to which they refer, but in some instances a more particular specification is made.

I take the opportunity of noticing a serious omission in my former Report, viz., I have not made mention of the elaborate memoir, Ostrogradsky ``M\'emoire sur les \'equations diff\'erentielles relatives au probl\`eme des Isop\'erim\`etres,'' M\'em.\ de St P\'et.,  t.\ iv.\ (6 s\'er.) pp.\ 385--517, 1850, which among other researches contains, and that in the most general form, the transformation of the equations of motion from the Lagrangian to the Hamiltonian form, and indeed the transformation of the general isoperimetric system (that is, the system arising from any problem in the calculus of variations) to the Hamiltonian form. I remark also, as regards the memoir of Cauchy referred to in the note p.\ 12 as an unpublished memoir of 1831, there is an ``Extrait du M\'emoire pr\'esent\'e \`a l'Acad\'emie de Turin le 11 Oct.\ 1831,'' published in lithograph under the date Turin, 1832, with an addition dated 6 Mar.\ 1833. The Extract begins thus:---``\S\ I. Variation des Constantes Arbitraires. Soient donn\'ees entre la variable $t,\ldots$ $n$ fonctions de $t$ d\'esign\'ees par $x,y,z\ldots$ et $n$ autres fonctions de $t$ d\'esign\'ees par $u,v,w,\ldots$ $2n$ \'equations diff\'erentielles du premier ordre et de la forme
\[ \frac{dx}{dt}=\frac{d\Omega}{du},\quad \frac{dy}{dt}=\frac{d\Omega}{dv},\quad \frac{dz}{dt}=\frac{d\Omega}{dw},\ldots \]
\[ \frac{du}{dt}=-\frac{d\Omega}{dx},\quad \frac{dv}{dt}=-\frac{d\Omega}{dy},\quad \frac{dw}{dt}=-\frac{d\Omega}{dz},\ldots'' \]
without any explanation as to the origin of these equations; and the formulae are then given for the variations of the constants in $\frac{dt}{-dz}$, the integrals of the foregoing system; this seems sufficient to establish that Cauchy in the year 1831 was familiar with the Hamiltonian form of the equations of motion.

% ============================================================
% PDF p.537 = printed p.515
% ============================================================
\newpage
\begin{center}
298] \hfill CERTAIN SPECIAL PROBLEMS OF DYNAMICS. \hfill 515
\end{center}

Bour's ``M\'emoire sur l'int\'egration des \'equations diff\'erentielles de la M\'ecanique,'' as published, M\'em.\ pr\'es.\ \`a l'Inst., t.\ xiv.\ pp.\ 792--821, is substantially the same as the extract thereof in Liouville's \emph{Journal}, referred to in my former Report; but since the date of that Report there have been published in the \emph{Comptes Rendus}, 1861 and 1862, several short papers by the same author; also Jacobi's great memoir, see list, Jacobi, \emph{Nova Methodus} \&c., 1862, edited after his decease by Clebsch; some valuable memoirs by Natani and Clebsch (\emph{Crelle}, 1861 and 1862) relating to the Pfaffian system of equations (which includes those of Dynamics), and Boole ``On Simultaneous Differential Equations of the First Order, in which the number of the Variables exceeds by more than one the number of the Equations,'' \emph{Phil.\ Trans.}, t.\ CLII.\ (1862), pp.\ 437--454.

\medskip
\begin{center}\textit{Rectilinear Motion. Article Nos.\ 1 to 5.}\end{center}

1. The determination of the motion of a falling body, which is the case of a constant force, is due to Galileo.

2. A variable force, assumed to be a force depending only on the position of the particle, may be considered as a function of the distance from any point in the line, selected at pleasure as a centre of force; but if, as usual, the force is given as a function of the distance from a certain point, it is natural to take that point for the centre of force. The problem thus becomes a particular case of that of central forces; and it is so treated in the \emph{Principia}, Book I.\ \S\ 7; the method has the advantage of explaining the paradoxical result which presents itself in the case Force $\propto(\text{Dist.})^{-2}$, and in some other cases where the force becomes infinite. According to theory, the velocity becomes infinite at the centre, but the direction of the motion is there abruptly reversed; so that the body in its motion does not pass through the centre, but on arriving there, forthwith returns towards its original position; of course such a motion cannot occur in nature, where neither a force nor a velocity ever is actually infinite.

3. Analytically the problem may be treated separately by means of the equation
\[ \frac{d^{2}x}{dt^{2}}=X, \]
which is at once integrable in the form $\left(\frac{dx}{dt}\right)^{2}=C+2\int X\,dx.$

4. The following cases may be mentioned:

Force $\propto$ Dist. The law of motion is well known, being in fact the same as for the cycloidal pendulum.

Force $\propto(\text{Dist.})^{-2}$, $=\frac{\mu}{x^{2}}$, which is the case above alluded to.

Assuming that the body falls from rest at a distance $a$, we have
\[ x = a(1-\cos\phi), \]
where, if $n=\frac{\sqrt{\mu}}{a^{\frac{3}{2}}}$, $\phi$ is given in terms of the time by means of the equation
\[ nt = \phi-\sin\phi. \]
\begin{flushright}\small 65---2\end{flushright}

% ============================================================
% PDF p.538 = printed p.516
% ============================================================
\newpage
\begin{center}
516 \hfill REPORT ON THE PROGRESS OF THE SOLUTION OF \hfill [298
\end{center}

If the body had initially a small transverse velocity, the motion would be in a very eccentric ellipse, and the formulae are in fact the limiting form of those for elliptic motion.

5. There are various laws of force for which the motion may be determined. In particular it can be determined by means of Elliptic Integrals, in the case of a body attracted to two centres, force $\propto(\text{dist.})^{-2}$: see Legendre, \emph{Exercices de Cal.\ Int\'eg.}, t.\ ii.\ pp.\ 502--512, and \emph{Th\'eorie des Fonct.\ Ellip.}, t.\ i.\ pp.\ 531--538.

\medskip
\begin{center}\textit{Central Forces, Article Nos.\ 6 to 26.}\end{center}

6. The theory of the motion of a body under the action of a given central force was first established in the \emph{Principia}, Book I.\ \S\S\ 2 and 3: viz.\ Prop.\ I.\ the areas are proportional to the times, that is (using the ordinary analytical notation), $r^{2}\,d\theta = h\,dt$, and Prop.\ VI.\ Cor.\ 3, $P\propto\frac{1}{SP^{2}\cdot PV}$, $=h^{2}u^{2}\left(\frac{d^{2}u}{d\theta^{2}}+u\right)$, so that
\[ \frac{d^{2}u}{d\theta^{2}}+u = \frac{P}{h^{2}u^{2}}. \]

7. It is to be noticed that, given the orbit, the law of force is at once determined; and \S\ 2 contains several instances of such determination; thus,

Prop.\ VII. If a body revolve in a circle, the law of force to a point $S$ is force $\propto\frac{1}{SP^{2}\cdot PV^{3}}$ ($P$ the body, $PV$ the chord through $S$).

Prop.\ IX. If a body move in a logarithmic spiral, force $\propto(\text{dist.})^{-3}$.

Prop.\ X. If a body move in an ellipse, force to centre $\propto$ dist., and as a particular case, if the body move in a parabola under the action of a force parallel to the axis, the force is constant. The particular case of motion in a parabola had been obtained by Galileo.

And \S\ 3, Props.\ XI.\ XII.\ XIII. If a body move in an ellipse, hyperbola, or parabola under the action of a force tending to the focus, force $\propto(\text{dist.})^{-2}$.

8. But Newton had no direct method of solving the inverse problem (which depends on the solution of the differential equation), ``Given the force to find the orbit.'' Thus force $\propto(\text{dist.})^{-2}$, after it has been shown that an ellipse, a hyperbola, and a parabola may each of them be described under the action of such a force, the remainder of the solution consists in showing that, given the initial circumstances of the motion, a conic section (ellipse, parabola, or hyperbola, as the case may be) can be constructed, passing through the point of projection, having its tangent in the direction of the initial motion, and such that the velocity of the body describing the conic section under the action of the given central force is equal to the velocity of projection; which being so, the orbit will be the conic section so constructed. This is what is done, Prop.\ XVII.; it may be observed that the latus rectum is constructed not very elegantly by means of the latus rectum of an auxiliary orbit.

% ============================================================
% PDF p.539 = printed p.517
% ============================================================
\newpage
\begin{center}
298] \hfill CERTAIN SPECIAL PROBLEMS OF DYNAMICS. \hfill 517
\end{center}

9. A more elegant construction was obtained by Cotes (see the \emph{Harmonia Mensurarum}, pp.\ 103--105, and demonstration from the author's papers in the Notes by R.\ Smith, pp.\ 124, 125), depending on the position of a point $C$, such that the velocity acquired in falling under the action of the central force from $C$ directly or through infinity\footnote{In the second case $C$ lies on the radius vector produced beyond the centre, and the body is supposed to fall from rest at $C$ (under the action of the central force considered as repulsive) to infinity, and then from the opposite infinity (with an initial velocity equal to the velocity so acquired) to $P$.} to $P$ the point of projection, is equal to the given velocity of projection.

10. But Newton's original construction is now usually replaced by a construction which employs the space due to the velocity of projection considered as produced by a constant force equal to the central force at the point of projection.

11. \S\ 9 of Book I.\ relates to revolving orbits, viz., it is shown that a body may be made to move in an orbit revolving round the centre of force, by adding to the central force required to make the body move in the same orbit at rest, a force $\propto(\text{dist.})^{-3}$. This appears very readily by means of the differential equation (\emph{ante}, No.\ 6), viz.\ writing therein $P+cu^{3}$ for $P$, and then $\theta'$, $h'$ in the place of $\theta\sqrt{1-\frac{c}{h^{2}}}$, $h\sqrt{1-\frac{c}{h^{2}}}$ respectively, the equation retains its original form, with $\theta'$, $h'$ in the place of $\theta$, $h$ respectively.

12. It may be remarked that when the original central force vanishes, the fixed orbit is a right line (not passing through the centre of force). It thus appears by \S\ 9 that the curve $u=A\cos(n\theta+B)$ may be described under the action of a force $\propto(\text{dist.})^{-3}$. A proposition in \S\ 2, already referred to, shows that a logarithmic spiral may be described under the action of such a force.

13. But the case of a force $\propto(\text{dist.})^{-3}$ was first completely discussed by Cotes in the \emph{Harmonia Mensurarum}, pp.\ 31--35, 98--104, and Notes, pp.\ 117--173. There are in all five cases, according as the velocity of projection is

\quad 1. Less than that acquired in falling from infinity, or say equal to that acquired in falling from a point $C$ to $P$, the point of projection.

\textit{[Figure: a triangle-like diagram with vertex $C'$ at top, point $P$ on the right, $C$ below $P$ near the bottom, and $S$ at lower-left; $Q$ is a point on segment from $P$ towards $C'$. This illustrates the geometric configuration referred to in the Cotes construction.]}

% ============================================================
% PDF p.540 = printed p.518
% ============================================================
\newpage
\begin{center}
518 \hfill REPORT ON THE PROGRESS OF THE SOLUTION OF \hfill [298
\end{center}

\quad 2. Equal to that acquired in falling from infinity.

\quad 3, 4, 5. Greater than that acquired in falling from infinity, or say equal to that acquired in falling from a point $C$, \emph{through} infinity, to $P$; viz.\ $PQ$ being the direction of projection, and $SQ$, $CT$ perpendiculars thereon from $S$ and $C$ respectively,
\begin{quote}
3. $SQ<TQ$;\\
4. $SQ=TQ$;\\
5. $SQ>TQ$;
\end{quote}
the equations of the orbits being

\quad 1. $u=Ae^{m\theta}+Be^{-m\theta}$, $A$ and $B$ same sign, so that rad.\ vector is never infinite.

\quad 2. $u=Ae^{m\theta}$ or $Be^{-m\theta}$, logarithmic spiral.

\quad 3. $u=Ae^{m\theta}+Be^{-m\theta}$, $A$ and $B$ opposite signs, so that rad.\ vector becomes infinite.

\quad 4. $u=A\theta+B$, $m=0$, reciprocal spiral.

\quad 5. $u=A\cos(m\theta+B)$, $m=n\sqrt{-1}$.

14. The before-mentioned equation,
\[ \frac{d^{2}u}{d\theta^{2}}+u = \frac{P}{h^{2}u^{2}}, \]
is in effect given (but the equation is encumbered with a tangential force) in Clairaut's \emph{Th\'eorie de la Lune}, 1765. It is given in its actual form, and extensively used (in particular for obtaining the above-mentioned equations for Cotes' spirals) in Whewell's \emph{Dynamics}, 1823. The equation appears to be but little known to continental writers, and (under the form $u''+u-aXu^{-3}=0$) it is given as new by Schellbach as late as 1853. The formulae used in place of it are those which give $t$ and $\theta$ each of them in terms of $r$; viz.
\[ dt = \frac{r\,dr}{\{-h^{2}+r^{2}(C-2\int P\,dr)\}^{\frac{1}{2}}}, \]
\[ d\theta = \frac{h\,dr}{r\{-h^{2}+r^{2}(C-2\int P\,dr)\}^{\frac{1}{2}}}, \]
which, however, assume that $P$ is a function of $r$ only.

15. Force $\propto(\text{dist.})^{-2}$. The law of motion in the conic sections is implicitly given by Newton's theorem for the equable description of the areas. For the parabola, if $a$ denote the pericentric distance, and $f$ the angle from pericentre or true anomaly, we have
\[ t = \frac{a^{\frac{3}{2}}}{\sqrt{\mu}}\left(\tan\tfrac{1}{2}f+\tfrac{1}{3}\tan^{3}\tfrac{1}{2}f\right). \]

% ============================================================
% PDF p.541 = printed p.519
% ============================================================
\newpage
\begin{center}
298] \hfill CERTAIN SPECIAL PROBLEMS OF DYNAMICS. \hfill 519
\end{center}

For the ellipse we have an angle $g$, the mean anomaly varying directly as the time $\left(g=nt\text{ if }n=\frac{\sqrt{\mu}}{a^{\frac{3}{2}}}\right)$; an auxiliary angle $u$, the eccentric anomaly, connected with $g$ by the equation
\[ g = u-e\sin u; \]
and then the radius vector $r$ and the true anomaly $f$ are given in terms of $u$ by the equations $r=a(1-e\cos u)$, and
\[ \cos f = \frac{\cos u-e}{1-e\cos u},\quad \sin f = \frac{\sqrt{1-e^{2}}\sin u}{1-e\cos u},\quad \text{and therefore}\quad \tan\tfrac{1}{2}f = \sqrt{\frac{1+e}{1-e}}\tan\tfrac{1}{2}u. \]

16. It is very convenient to have a notation for $\frac{r}{a}$ and $f$ considered as functions of $e$, $g$, and I have elsewhere proposed to write
\[ r = a\,\mathrm{elqr}(e,g),\quad f = \mathrm{elta}(e,g), \]
read elqr elliptic quotient radius, and elta elliptic true anomaly.

17. The formulae for the hyperbola correspond to those for the ellipse, but they contain exponential in the place of circular functions (see post, Elliptic Motion).

18. Euler, in the memoir ``Determinatio Orbitae Cometae Anni 1742,'' (1743), p.\ 16 \emph{et seq.}, obtained an expression for the time of describing a parabolic arc in terms of the radius vectors and the chord; viz.\ these being $f$, $g$, and $k$, the expression is
\[ t\sqrt{\mu} = \tfrac{1}{6}\sqrt{2}\{(f+g+k)^{\frac{3}{2}}-(f+g-k)^{\frac{3}{2}}\}, \]
which, however, as remarked by Lagrange, \emph{M\'ec.\ Anal.}, t.\ ii.\ (3rd edit.\ p.\ 28), is deducible from Lemma X.\ of the third book of the \emph{Principia}. But the theorem in its actual form is due to Euler.

19. Lambert, in the \emph{Proprietates Insigniores} \&c.\ (1761), Theorem VII.\ Cor.\ 2, obtained the same theorem, and in section 4 he obtained the corresponding theorem for elliptic motion; viz.\ the expression for the time is
\[ t = \frac{a^{\frac{3}{2}}}{\sqrt{\mu}}\{\phi-\phi'-(\sin\phi-\sin\phi')\}, \]
if
\[ \sin\tfrac{1}{2}\phi = \tfrac{1}{2}\sqrt{\frac{f+g+k}{a}},\quad \sin\tfrac{1}{2}\phi' = \tfrac{1}{2}\sqrt{\frac{f+g-k}{a}}. \]
The form of the formula is, it will be observed, similar to that for motion in a straight line (\emph{ante}, No.\ 4), and in fact the motion in the ellipse is, by an ingenious geometrical transformation, made to depend upon that in the straight line. The geometrical theorems upon which the transformation depends are stated, Cayley ``On Lambert's Theorem \&c.'' (1861).

% ============================================================
% PDF p.542 = printed p.520
% ============================================================
\newpage
\begin{center}
520 \hfill REPORT ON THE PROGRESS OF THE SOLUTION OF \hfill [298
\end{center}

20. The theorem was also obtained by Lagrange in the memoir ``Recherches \&c.'' (1767) as a corollary to his solution of the problem of two centres; viz.\ upon making the attractive force of one of the centres equal to zero, and assuming that such centre is situate on the curve, the expression for the time presents itself in the form given by Lambert's theorem.

21. Two other demonstrations of the theorem are given by Lagrange in the memoir ``Sur une mani\`ere particuli\`ere d'exprimer le temps \&c.'' (1778), reproduced in Note V.\ of the second volume of the last edition (Bertrand's) of the \emph{M\'ecanique Analytique}. As M.\ Bertrand remarks, these demonstrations are very complete, very elegant, and very natural, assuming that the theorem is known beforehand.

Demonstrations were also given by Gauss, ``Theoria Motus'' (1809), p.\ 119 \emph{et seq.}; Pagani, ``D\'emonstration d'un th\'eor\`eme \&c.'' (1834); and (in connexion with Hamilton's Principal Function) by Sir W.\ R.\ Hamilton, ``On a General Method \&c.'' (1834), p.\ 282; Jacobi, ``Zur Theorie \&c.'' (1837), p.\ 122; Cayley, ``Note on the Theory of Elliptic Motion'' (1856).

22. Connected with the problem of central forces, we have Sir W.\ R.\ Hamilton's ``Hodograph,'' which in the paper (\emph{Proc.\ R.\ Irish Acad.}\ 1847) is defined, and the fundamental properties thereof are stated; viz.\ if in an orbit round a centre of force there be taken on the perpendicular from the centre on the tangent at each point, a length equal to the velocity at that point of the orbit, the extremities of these lengths will trace out a curve which is the hodograph. As the product of the velocity into the perpendicular on the tangent is equal to twice the area swept out in a unit of time ($vp=h$), the hodograph is the reciprocal polar of the orbit with respect to a circle described about the centre of force, radius $=\sqrt{h}$. Whence also the tangent at any point of the hodograph is perpendicular to the radius vector through the corresponding point of the orbit, and the product of the perpendicular on the tangent into the corresponding radius vector is $=h$.

If force $\propto(\text{dist.})^{-2}$, the hodograph, \emph{qu\^a} reciprocal polar of a conic section with respect to a circle described about the focus, is a circle.

23. The following theorem is also given without demonstration; viz.\ if two circular hodographs, which have a common chord passing or tending through a common centre of force, be both cut at right angles by a third circle, the times of hodographically describing the intercepted arcs (that is, the times of describing the corresponding elliptic arcs) will be equal.

24. Droop, ``On the Isochronism \&c.'' (1856), shows geometrically that the last-mentioned property is equivalent to Lambert's theorem; and an analytical demonstration is also given, Cayley, ``A demonstration of Sir W.\ R.\ Hamilton's Theorem \&c.'' (1857). See also Sir W.\ R.\ Hamilton's \emph{Lectures on Quaternions} (1853), p.\ 614.

25. The laws of central force which have been thus far referred to are force $\propto r$, $\propto\frac{1}{r^{2}}$, $\propto\frac{1}{r^{3}}$; and it has been seen that the case of a force $P+\frac{C}{r^{3}}$ depends upon that

% ============================================================
% PDF p.543 = printed p.521
% ============================================================
\newpage
\begin{center}
298] \hfill CERTAIN SPECIAL PROBLEMS OF DYNAMICS. \hfill 521
\end{center}

of a force $P$, so that the motions for the forces $Ar+\frac{C}{r^{3}}$ and $\frac{B}{r^{2}}+\frac{C}{r^{3}}$ are deducible from those for the forces $Ar$ and $\frac{B}{r^{2}}$ respectively. Some other laws of force, e.g.\ $\frac{A}{r^{2}}\pm Br$, $\frac{A}{r^{2}}+\frac{B}{r^{3}}+\frac{C}{r^{4}}+\frac{D}{r^{5}}$, are considered by Legendre, ``Th\'eorie des Fonctions Elliptiques'' (1825), being such as lead to results expressible by elliptic integrals, and also the law $\frac{M}{r^{2}}$, for which the result involves a peculiar logarithmic integral. But the most elaborate examination of the different cases in which the solution can be worked out by elliptic integrals or otherwise is given in Stader's memoir ``De Orbitis \&c.'' (1852), where the investigation is conducted by means of the formulae which give $t$ and $\theta$ in terms of $r$ (\emph{ante}, No.\ 14).

26. In speaking of a central force, it is for the most part implied that the force is a function of the distance: for some problems in which this is not the case, see post, Miscellaneous Problems, Nos.\ 86 and 87.

It is to be noticed that, although the problem of central forces may be (as it has so far been) considered as a problem \emph{in plano} (viz.\ the plane of the motion has been made the plane of reference), yet that it is also interesting to consider it as a problem in space; in fact, in this case the integrals, though of course involved in those which belong to the plane problem, present themselves under very distinct forms, and afford interesting applications of the theory of canonical integrals, of the derivation of the successive integrals by Poisson's method, and of other general dynamical theories. Moreover, in the lunar and planetary theories, the problem must of necessity be so treated. Without going into any details on this point, I will refer to Bertrand's memoir, ``Sur les \'equations diff\'erentielles de la M\'ecanique'' (1852), Donkin's memoir ``On a Class of Differential Equations \&c.'' (1855), and Jacobi's posthumous memoir, ``Nova Methodus \&c.'' (1862).

\medskip
\begin{center}\textit{Elliptic Motion, Article Nos.\ 27---40.}\end{center}

27. The question of the development of the true anomaly in terms of the mean anomaly (Kepler's problem), and of the other developments which present themselves in the theory of elliptic motion, is one that has very much occupied the attention of geometers. The formulae on which it depends are mentioned \emph{ante}, No.\ 15; they involve as an auxiliary quantity the eccentric anomaly $u$.

28. Consider first the equation
\[ g = u-e\sin u, \]
which connects the mean anomaly $g$ with the eccentric anomaly $u$.

Any function of $u$, and in particular $u$ itself, and the functions $\genfrac{}{}{0pt}{}{\cos}{\sin}\,nu$ may be expanded in terms of $g$ by means of Lagrange's theorem (Lagrange, \emph{M\'em.\ de Berlin}, 1768--1769, ``Th\'eorie des Fonctions,'' chap.\ 16, and ``Trait\'e de la R\'esolution des \'equations Num\'eriques,'' Note 11).

\vfill
\begin{flushright}\small C.\ IV.\quad 66\end{flushright}

% ============================================================
% PDF p.544 = printed p.522
% ============================================================
\newpage
\begin{center}
522 \hfill REPORT ON THE PROGRESS OF THE SOLUTION OF \hfill [298
\end{center}

29. Considering next the equation
\[ \tan\tfrac{1}{2}f = \sqrt{\frac{1+e}{1-e}}\tan\tfrac{1}{2}u, \]
which gives the true anomaly in terms of the eccentric anomaly, then, by replacing the circular functions by their exponential values (a process employed by Lagrange, \emph{M\'em.\ de Berlin}, 1776), $f$ can be expressed in terms of $u$; viz.\ the result is
\[ f = u+2\lambda\sin u+2\lambda^{2}\cdot\tfrac{1}{2}\sin 2u+2\lambda^{3}\cdot\tfrac{1}{3}\sin 3u+\&\text{c.}, \]
where $\lambda = \frac{1-\sqrt{1-e^{2}}}{e}\left(=\frac{e}{1+\sqrt{1-e^{2}}}\right)$. Hence if $u$, $\sin u$, $\sin 2u$, \&c.\ are expressed in terms of the mean anomaly, $f$ will be obtained in the form $f=g+$ a series of multiple sines of $g$, the coefficients of the different terms being given in the first instance as functions of $e$ and $\lambda$; and to complete the development $\lambda$ and its powers have to be developed in powers of $e$. The solution is carried thus far in the \emph{M\'ecanique Analytique} (1788), and in the \emph{M\'ecanique C\'eleste} (1799).

30. We have next Bessel's investigations in the Berlin Memoirs for 1816, 1818, and 1824, and which are carried on mainly by means of the integral
\[ 2\pi I_{h}^{k} = \int_{0}^{2\pi}\cos(hu-k\sin u)\,du, \]
and various properties are there obtained and applications made of this important transcendant.

31. Relating to this integral we have Jacobi's memoir, ``Formulae transformationis \&c.'' (1836), Liouville, ``Sur l'int\'egrale $\int_{0}^{2\pi}\cos i(u-x\sin u)\,du$'' (1841), and Hansen's ``Ermittelung der absoluten St\"orungen'' (1843); the researches of Poisson in the \emph{Connaissance des Temps} for 1825 and 1836 are closely connected with those of Bessel.

32. A very elegant formula, giving the actual expression of the coefficients considered as functions of $e$ and $\lambda$, is given by Greatheed in the paper ``Investigation of the General Term \&c.'' (1838); viz.\ this is
\[ f = g+2\Sigma\lambda^{r}\,\frac{e^{\frac{1}{2}r(\lambda+\lambda^{-1})}+\lambda^{-r}e^{-\frac{1}{2}r(\lambda+\lambda^{-1})}}{r}\sin rg, \]
where, after developing in powers of $\lambda$, the negative powers of $\lambda$ must be rejected, and the term independent of $\lambda$ divided by 2. This result is extended to other functions of $f$, Cayley ``On certain Expansions \&c.'' (1842).

33. An expression for the coefficient of the general term as a function of $e$ only is obtained, Lefort, ``Expression Num\'erique \&c.'' (1846). The expression, which, from the nature of the case, is a very complicated one, is obtained by means of Bessel's integral. This is an indirect process which really comes to the combination of the developments of $f$ in terms of $u$, and $u$ in terms of $g$; and an equivalent result is obtained directly in this manner, Greedy, ``General and Practical Solution \&c.'' (1855).

% ============================================================
% PDF p.545 = printed p.523
% ============================================================
\newpage
\begin{center}
298] \hfill CERTAIN SPECIAL PROBLEMS OF DYNAMICS. \hfill 523
\end{center}

34. We have also on the subject of these developments the very valuable and interesting researches of Hansen, contained in his \emph{Fundamenta Nova} \&c.\ (1838), in the memoir ``Ermittelung der absoluten St\"orungen \&c.'' (1845), and in particular in the memoir ``Entwickelung des Products \&c.'' (1853).

35. But the expression for the coefficient of the general term $\genfrac{}{}{0pt}{}{\cos}{\sin}\,rg$ in any of these expansions is so complicated that it was desirable to have for the coefficients corresponding to the values $r=0,1,2,3,\ldots$ the finally reduced expressions in which the coefficient of each power of $e$ is given as a numerical fraction. Such formulae for $\left(\frac{r}{a}\right)^{n}\genfrac{}{}{0pt}{}{\cos}{\sin}\,jf$, where $j$ is a general symbol, the expansion being carried as far as $e^{6}$, were given, Leverrier, \emph{Annales de l'Observatoire de Paris}, t.\ i.\ (1855).

36. And starting from these I deduced the results given in my ``Tables of the Developments \&c.'' (1861); viz.\ these tables give $\left(x=\frac{r}{a}\right)$,
\[ \left(\frac{r}{a}\right)^{n}\genfrac{}{}{0pt}{}{\cos}{\sin}\,jf,\quad (n=-5,\ldots,5;\ j=1,\ldots,7), \]
all carried to $e^{7}$.

37. The true anomaly $f$ has been repeatedly calculated to a much greater extent, in particular by Schubert (\emph{Ast.\ Th\'eorique}, St Pet.\ 1822), as far as $e^{20}$. The expression for $\frac{r}{a}$ as far as $e^{20}$ is given in the same work, and that for $\log\frac{r}{a}$ as far as $e^{20}$ was calculated by Oriani, see Introd.\ to Delambre's \emph{Tables du Soleil}, Paris (1806).

38. It may be remarked that when the motion of a body is referred to a plane which is not the plane of the elliptic orbit, then we have questions of development similar in some measure to those which regard the motion in the orbit; if, for instance, $z$ be the distance from node, $\phi$ the inclination, and $x$ the reduced distance from node, then $\cos\phi = \cos z\cos x$, from which we may derive $z=x+$ series of multiple sines of $x$. And there are, moreover, the questions connected with the development of the reciprocal distance of two particles---say $(a^{2}+a'^{2}-2aa'\cos\theta)^{-\frac{1}{2}}$---which present themselves in the planetary theory; but this last is a wide subject, which I do not here enter upon. I will, however, just refer to Hansen's memoir, ``Ueber die Entwickelung der negativen und ungeraden Potenzen \&c.'' (1854).

\begin{flushright}\small 66---2\end{flushright}

\end{document}
